% Native qualitative models; selected source units remain version-distinct.
\chapter{Qualitative Processes and Compositional Models}
\label{sec:p16-qualitative-process-and-model-reasoning}

Qualitative reasoning represents distinctions which matter to a task without
requiring every numerical value. A rising level, an exhausted supply, a
pressure difference and an active heat flow are different kinds of information.
They must be represented with constraints which say how they interact, not
merely described as qualitative. Forbus's 2008 account makes discretization,
relevance and ambiguity joint design principles
\cite[secs.~9.1--9.6]{Forbus2008QualitativeModeling}.

Three source regimes are distinguished here. Qualitative process theory (QPT)
organizes changes through processes and individual views. De Kleer and Brown's
component framework uses mode-dependent confluences and an envisionment.
Kuipers's 1986 QSIM operates on an explicit qualitative differential equation,
ordered landmarks and point/interval transitions. A translation between their
notations does not erase their different closure, time and function assumptions.
The quantitative structural models of \cref{def:p16-finite-scm} do not
substitute for these representations or their model-formulation methods.

\section{Signs, quantities and relevant comparisons}
\label{sec:p16-qr-quantities}

\begin{hbdefinition}{def:p16-qr-signs}{A precise sign abstraction}
For real \(x\), let \(\sigma(x)\in\{-1,0,1\}\) be its sign.
For signs \(s,t\), define \(s\oplus t\) to be the set of signs
\(\sigma(x+y)\) obtainable with \(\sigma(x)=s,\sigma(y)=t\).
Define \(s\otimes t=\sigma(xy)\) under those same restrictions.
The sum is set-valued: an ambiguous result is a set of possibilities,
not a fourth measured numerical value.
\end{hbdefinition}

\begin{hbtheorem}{thm:p16-qr-signs}{proved=annex:p16.qr-signs}{Exact local sign operations}
The additive table is
\[
\begin{array}{c|ccc}
 \oplus&-1&0&1\\\hline
 -1&\{-1\}&\{-1\}&\{-1,0,1\}\\
 0&\{-1\}&\{0\}&\{1\}\\
 1&\{-1,0,1\}&\{1\}&\{1\}
\end{array}
\]
and multiplication has the ordinary product of signs, with zero absorbing.
Every concrete sum/product has an admitted abstract result, and every result
in an individual table entry has a concrete witness. This is a two-input
statement, not global realizability of a network or simulation.
\end{hbtheorem}

De Kleer and Brown use qualitative values and confluences to propagate
constraints within a component mode. Opposing tendencies cannot be cancelled
as if their magnitudes were known. Repeated occurrences of the same quantity
also carry an identity which independent table lookups can lose. For example,
\(x-x=0\) exactly; treating its two signed terms as independent positive
and negative inputs would admit all three results.

QPT's native quantity has an amount \(A(q)\) and derivative \(D(q)\).
Their signs and magnitudes are further distinguished, including the derivative
sign \(D_s(q)\). A quantity space records ordinal comparisons with relevant
other values, and can be only partially ordered. A \emph{completion} supplies
the missing order comparisons; it is a considered alternative, not knowledge
that was already available \cite[secs.~2.3--2.4]{Forbus1984QPT}.

Forbus's 2008 account distinguishes a fixed numerical landmark from a more
general limit point. A neighboring object's changing temperature can be a
limit point relevant to heat flow without being a fixed landmark. A totally
ordered value space with \(n\) distinct finite limit points has \(n\)
point values and \(n+1\) intervening/exterior interval values. Introducing
new landmarks during reasoning increases resolution but can make the state
space unbounded \cite[pp.~367--368]{Forbus2008QualitativeModeling}.

\section{Processes, activity and influence resolution}
\label{sec:p16-qr-processes}

QPT represents a process by its participating individuals, external
preconditions, quantity conditions, relations and direct influences.
Matching the participant description creates a possible process instance.
The instance is active only when its preconditions and quantity conditions
hold. Relations and influences are then enabled. An individual view has
the related role of describing contingent objects or modes, but does not
itself supply direct influences \cite[secs.~2.5, 3.1--3.6]{Forbus1984QPT}.

\begin{hbexample}{def:p16-qr-heat-process}{The native heat-flow organization}
The heat-flow participants are a source and destination with heat quantities,
and a connecting heat path. Heat alignment is an external precondition;
the source being hotter than the destination is a quantity condition.
When active, a positive flow rate depends monotonically on the temperature
difference. The process supplies \(I^-(\mathit{heat}(src),A(rate))\)
and \(I^+(\mathit{heat}(dst),A(rate))\).

A sealed path can still be a potential connection without an active flow.
Likewise the existence of a contained liquid and the amount remaining in it
are separate from whether a process affecting it is active. Closing a valve
is an external action, not something inferred merely from a derivative sign.
\end{hbexample}

The original QPT grammar separates preconditions from dynamically predictable
quantity conditions. It also imposes a sole-mechanism assumption, supplemented
by completeness assumptions about the process vocabulary and dependencies
when reasoning by exclusion. In its domain model a quantity is directly
influenced, indirectly influenced, or uninfluenced, not both directly and
indirectly influenced at once. The active qualitative-proportionality graph
is loop-free; feedback can still run through a derivative/direct-influence
relationship \cite[secs.~3.2--3.6]{Forbus1984QPT}.

\begin{hbdefinition}{def:p16-qr-balance}{A fully stated local heat balance}
On a time interval, let \(E\) be a differentiable heat-stock parameter.
An active heating process supplies positive rate \(u\), an active loss
process supplies positive rate \(v\), and these are all its direct influences.
Their additive influence semantics gives \(E'=u-v\). The loss process's
own activation condition, a temperature difference and an open path, remains
true in the scenario being described.
\end{hbdefinition}

\begin{hbtheorem}{thm:p16-qr-balance}{proved=annex:p16.qr-balance}{A stationary stock need not mean inactive processes}
At each time, \(D_s(E)\) is positive, zero or negative according as
\(u>v\), \(u=v\) or \(u<v\). Knowing only that both rates are
positive leaves all three cases possible. When the rates are equal, both
processes can remain active while their net stock derivative is zero.
\end{hbtheorem}

Direct influences combine additively; indirect qualitative proportionalities
describe partial functional dependencies and need not combine that way.
For example, two increasing inputs with opposite effects on an output do
not determine which dominates. A closed-world assumption about the active
dependency set is consequential: adding another process or input can change
the inferred direction. QPT's causal commitments are not statistical
identification assumptions or a generic consequence of an equation's syntax.

\section{Monotonic dependencies are not interchangeable}
\label{sec:p16-qr-monotonicity}

In QPT, \(q_1\propto_Q^+q_2\) describes an increasing functional
dependence on \(q_2\) with potentially other inputs. Its all-else-equal
qualification must be supplied by the current process structure and modeling
assumptions. Kuipers's 1986 \(M^+(f,g)\) instead means
\(f(t)=H(g(t))\), with \(H'>0\) in the interior of its input range;
\(M^-\) uses \(H'<0\). This is stronger than merely calling a
function strictly increasing \cite[Definitions~3.5--3.6,
pp.~302--303]{Kuipers1986QualitativeSimulation}.

\begin{hbtheorem}{thm:p16-qr-monotone}{proved=annex:p16.qr-monotone}{The finite differentiable monotonic-constraint argument}
Let real-valued \(f,g\) be continuously differentiable on a compact
interval and \(f=H\circ g\), where \(H\) is continuous and strictly
increasing on the range of \(g\), and differentiable with positive
derivative in that range's interior. At an interior time, the derivative
signs of \(f,g\) agree. A strictly decreasing \(H\) with negative
interior derivative reverses the nonzero signs and preserves zero.

Strict increasing monotonicity alone does not imply this pointwise sign
identity: \(g(t)=t,H(x)=x^3\) gives \(g'(0)=1\) but \(f'(0)=0\).
Nor does a dependency with a changing additional input qualify as a
one-input \(M^+\) constraint.
\end{hbtheorem}

This proves the stated finite differentiable portion of the native constraint,
not QSIM's complete extended-endpoint coverage theorem. The direction
convention is also explicit: the first argument here is the dependent output,
as in the 1986 definition, not an unannounced interchange with earlier notation.
Corresponding landmark values add equality information; monotonicity alone
does not say that equal inputs in two different unknown functions produce
equal outputs. A shared function specification is another premise.

\section{Components and fields}
\label{sec:p16-qr-ontologies}

The component ontology represents terminal-connected devices with generic
component models. De Kleer and Brown's no-function-in-structure principle
requires those models not to encode the desired functioning of the whole.
A closed switch does not guarantee current: other connections and the
potential across it matter. The source qualifies context independence by
explicit class-wide assumptions, including its quasistatic approximation
\cite[secs.~2.3--2.5]{DeKleerBrown1984Confluences}.

For their idealized valve, WORKING, OPEN and CLOSED modes have different
constraints. In WORKING mode, changes in pressure \(P\), aperture \(A\)
and flow \(Q\) obey the sign confluence corresponding to
\(dP+dA-dQ=0\). OPEN is an idealized maximal-aperture mode; CLOSED
forces zero flow. These mode assumptions are not a detailed valve equation
valid at every frequency, pressure or geometry.

\begin{hbtheorem}{thm:p16-qr-valve}{proved=annex:p16.qr-valve}{A confluence inference with its scope exposed}
In the stated WORKING confluence, an increasing aperture and constant
flow force a decreasing pressure change. An increasing pressure with a
decreasing aperture leaves the flow change ambiguous without more information.
The variables in this result are changes at the stated qualitative level,
not interchangeable stock quantities or rate magnitudes.
\end{hbtheorem}

The source pressure regulator senses increased output pressure and moves a
diaphragm to close the valve, opposing the initial increase in flow.
Its coarse analysis captures attenuation, not an assertion that output
pressure never changes. Intrastate behavior solves the confluences for a
current mode; interstate reasoning checks whether a mode boundary can cease
to hold. ENVISION combines constraint propagation with generate-and-test
because simultaneous qualitative constraints cannot generally be solved by
ordinary field-algebra elimination \cite[secs.~4.1--4.2]{DeKleerBrown1984Confluences}.

Field ontologies, by contrast, retain spatially distributed structure.
Forbus describes regions with qualitatively equivalent field properties and
spatial aggregation from lower-level structures, such as isobar segments,
to higher-level patterns such as troughs. A single lumped pressure value
does not preserve those distinctions. Physical space, phase space and
configuration space are different domains for such partitions
\cite[sec.~9.3.3]{Forbus2008QualitativeModeling}.

\begin{hbexample}{def:p16-qr-field-boundary}{A spatial pattern lost by lumping}
Suppose three neighboring sampled regions have pressure categories
\(high,low,high\). Their central depression is different from
\(low,high,low\), although the same categories occur in both. Identifying
the middle region, its neighbors and their ordering is essential to the
pattern. This illustrates why region relations belong in the representation;
it is not a fluid-dynamics prediction from three labels. Contour construction,
geometric validity and evolution of the field require their own models.
\end{hbexample}

Process and component models are lumped-parameter organizations; field
models are not made equivalent to them merely by drawing more nodes.
Forbus explicitly notes that field-model qualitative simulation has a less
developed basis than the static analysis frameworks he surveys.
\section{Compositional model formulation}
\label{sec:p16-qr-model-formulation}

A scenario description, domain theory and query are distinct inputs.
Model fragments specify participants, applicability conditions, quantities and
consequences. Assumption classes organize alternative modeling choices;
conditional choices can make further classes relevant. Coherence requires
compatible choices, while adequacy concerns the specified query and accuracy
criterion. Having equations is not itself either property
\cite[sec.~9.5]{Forbus2008QualitativeModeling}.

\begin{hbexample}{def:p16-qr-formulation-example}{A fluid path at two levels of detail}
The source's fluid-path fragment makes conductance increase with path size
and decrease with path length, when the path is unblocked. A geometric
submodel may estimate size from width and height if the opening has an
approximately rectangular projection. Choosing an abstract path versus a
geometric path is one assumption class; selecting rectangular, circular or
irregular geometry becomes relevant only in the latter choice.

A question about whether flow is possible can often use an ordinal pressure
comparison. A question about its timing needs more information about rates
and the path. Treating a system as a black box while also instantiating its
subsystems can violate a declared coherence constraint. Likewise, ignoring
the working fluid's properties while retaining a boiler relief valve is not
automatically a coherent reduction. These are modeling restrictions, not
consequences of choosing the fewest symbols.
\end{hbexample}

\begin{hbmethod}{def:p16-qr-two-pass}{The source's two-pass formulation organization}
First instantiate potentially relevant model fragments against the structural
description. Record the assumptions supporting each consequence and nogoods
which exclude inconsistent combinations; an assumption-based truth-maintenance
system retains alternative supporting environments. Find coherent support
sets which constrain the query terms. Then discard the expansive preliminary
construction and reinstate only the selected relevant fragment instances.
The source's initial method uses smallest support sets as its simplicity
criterion; that criterion is a modeling premise, not physical truth.
\end{hbmethod}

Forbus contrasts this with restricted searches using independent assumption
classes and partially ordered simplicity choices, and with query-directed
time-scale selection. Fast phenomena can sometimes be replaced by functional
relations and slow ones ignored, but those approximations are query-relative.
The independence and adequacy restrictions must not be dropped while retaining
an efficiency claim. No polynomial-time assertion for unrestricted model
formulation follows from this account.

\section{A native multi-process state and its limits}
\label{sec:p16-qr-containers}

\begin{hbdefinition}{def:p16-qr-containers}{The three-container comparison at explicit scope}
Following the source scenario, let \(F,G,H\) be open containers with
homogeneous liquid, common bottom elevation, open connecting paths \(G-F\)
and \(G-H\), and no other mass flows. Let \(V_i\) be their amounts and
\(h_i\) their heads, with differentiable increasing head functions whose
derivatives \(\beta_i>0\) in the current mode. Pressure comparisons use
the same hydrostatic relation. Initially \(h_G>h_F,h_H\).
Let \(q_F,q_H\) be signed net rates oriented from \(G\) toward the
other container: they have the sign of the corresponding head difference,
and are zero at equal heads. Initially the two active process instances
therefore have positive rates, with direct influences
\[
 V_F'=q_F,\quad V_H'=q_H,\quad V_G'=-q_F-q_H.
\]
The formulation supplies stock/rate types and a complete influence set;
it does not assign their unknown magnitudes.
\end{hbdefinition}

\begin{hbtheorem}{thm:p16-qr-containers}{proved=annex:p16.qr-containers}{Resolved directions, conserved stock and a boundary reversal}
In the initial mode, \(G\)'s amount/head decrease and the other two increase.
Both head differences \(h_G-h_F,h_G-h_H\) decrease, and total amount is
constant. Their derivative signs alone do not determine which difference
reaches zero first, or whether equality is simultaneous.

If \(h_G=h_F>h_H\) is reached, the \(G-F\) rate is zero, while
the \(G-H\) rate remains positive. At that boundary the derivative of
\(h_G-h_F\) is negative. With a right-continuous derivative and retained
open path, the head order immediately reverses locally, enabling \(F-G\)
flow. The symmetric case exchanges \(F,H\). This is conditional on the
boundary actually being reached and on the stated mode assumptions.
\end{hbtheorem}

Forbus's Figure~9.4 distinguishes the two possible first equalizations and
the simultaneous alternative. Its instant-labeled transitions represent
departures from equality; they are not measured travel times. A raised
container bottom or trapped gas can reverse a pressure comparison based
only on local liquid levels. The initial-model assumptions therefore matter
\cite[pp.~363--365, Figure~9.4]{Forbus2008QualitativeModeling}.

\begin{hbtheorem}{thm:p16-qr-contact}{proved=annex:p16.qr-contact}{Continuity, point stationarity and finite contact}
Continuous \(x,l\) cannot change \(x-l\) from negative to positive
without an equality time. Positive derivative while below a fixed landmark
does not alone force finite contact: \(x(t)=l-e^{-t}\) never reaches it
at a finite \(t\). A uniform lower bound \((x-l)'\ge\epsilon>0\),
retained until contact over an interval long enough, does force contact within
\((l(t_0)-x(t_0))/\epsilon\). Zero derivative at one point is not an
interval of constant value, as \(x(t)=t^2\) at zero demonstrates.
For the derivative/integration assertions assume \(x,l\) are continuously
differentiable throughout the stated comparison interval.
\end{hbtheorem}

QPT's limit analysis compares changing quantities with neighboring limit
points, tests single and simultaneous inequality changes, and revises process
and view activity. Its original finite-reaching and equality-change assumptions
must be stated, not obtained from a sign alone
\cite[sec.~3.6.4 and footnote~13]{Forbus1984QPT}. QSIM instead explicitly
uses standard real time, distinguished points and open intervals, and permits
infinite endpoints for asymptotic limits. These regimes are not interchangeable.

\section{QSIM's actual constraint workflow}
\label{sec:p16-qr-qsim}

Kuipers's 1986 state pairs a landmark/interval magnitude with a direction
\(inc,std,dec\). Its structure uses ADD, MULT, MINUS, DERIV and strong
\(M^+,M^-\) constraints, corresponding landmark values, operating-region
limits and an initial state. DERIV\((f,g)\) means \(f'=g\), not that
the signs of \(f,g\) are equal. A new extremum may create a new landmark
\cite[secs.~2--4]{Kuipers1986QualitativeSimulation}.

\begin{hbmethod}{def:p16-qr-qsim-method}{The six-stage 1986 organization}
Select a state from ACTIVE. Generate each parameter's legal point-to-interval
or interval-to-point transitions. Construct transition tuples for individual
constraints and filter using direction and correspondence restrictions.
Apply pairwise consistency between constraints sharing parameters. Join
the surviving tuples into global assignments, rejecting conflicting choices.
Construct successor states and apply the pure no-change, exact-cycle and
divergence filters before returning appropriate successors to ACTIVE.
\end{hbmethod}

The algorithm filters transitions rather than merely states. Two neighboring
constraints must assign the same transition to their shared parameter.
Global joining is still necessary: in the source's two adjacent monotonic
constraints, the alternatives \((I2,I2)\) and \((I3,I3)\) yield only
two global choices, not four independent combinations. Heuristic filters,
including quiescence or excluded divergence, are distinct from the pure
algorithm and require their own warrants
\cite[secs.~4.5--4.9]{Kuipers1986QualitativeSimulation}.

\begin{hbtheorem}{thm:p16-qr-filter}{proved=annex:p16.qr-filter}{What finite tuple pruning actually preserves}
For finite transition domains and finite constraint-tuple sets, repeatedly
delete a tuple without matching support in a neighboring constraint on their
shared variables. Deletion terminates and preserves every already satisfying
global assignment. It need not produce one: binary variables with pairwise
inequality constraints on a triangle retain local supports but have no global
solution. This is a finite constraint lemma, not full QSIM trajectory coverage.
\end{hbtheorem}

The source's upward-thrown-ball example uses DERIV\((Y,V)\),
DERIV\((V,A)\) and constant negative \(A\). The apex has \(V=0\)
and hence \(Y'=0\); its height is a new landmark, not a numerical altitude.
The directions and correspondences eliminate inconsistent local transition
tuples. A qualitative plot constant in its drawn interval label does not mean
the underlying quantity is constant.

\section{Why local admissibility is not realizability}
\label{sec:p16-qr-spurious}

Kuipers's frictionless-spring example can branch at a return comparison
according to velocity reaching its earlier landmark before, with, or after
position returns to zero. The energy invariant distinguishes the physically
admissible return from the extra qualitative alternatives
\cite[sec.~5.3]{Kuipers1986QualitativeSimulation}.

\begin{hbtheorem}{thm:p16-qr-energy}{proved=annex:p16.qr-energy}{The full invariant argument for the stated restoring model}
Let \(x'=v,v'=-H(x)\), with real continuously differentiable solution
on an interval, continuous restoring function \(H\), and \(H(0)=0\).
Then
\[
 {\cal E}(t)=\tfrac12v(t)^2+\int_0^{x(t)}H(y)\,dy
\]
is constant. If \(x(t_0)=x(t_1)=0\) and both velocities are positive,
then \(v(t_1)=v(t_0)\). Unequal positive-velocity returns to zero are
not realizations of this model. The result assumes no damping or external work.
\end{hbtheorem}

Increasing, steady and decaying amplitudes cannot be distinguished merely
by repeating the same signs. New extrema landmarks make the question
expressible, but expressing it is not deciding it. An energy constraint can
rule out spurious paths for this system; automatically discovering such
global invariants for arbitrary qualitative models is a different problem.
Neither the local sign table nor the pruning lemma establishes QSIM's full
guaranteed-coverage theorem, its complete correspondence filters or their
extended-endpoint cases.

\begin{hbmethod}{def:p16-qr-consultation}{Using a qualitative model as a reference method}
State the task, quantity types, comparison points and operating modes.
Identify actual participant/view/process or component instances, external
conditions and class-wide approximations. Record closure assumptions and
all active effects before resolving changes. Preserve alternative orders,
unknown relative rates and unresolved next-event timing. Construct and
check model choices before simulating; retain the distinction between local
consistency, global constraint satisfaction and a realizable trajectory.
Seek finer models or observations for ambiguous or potentially spurious
behaviors. Do not convert a generated hazard into a prediction that it must occur.
\end{hbmethod}

\index{qualitative process theory}\index{quantity space}\index{confluence}
\index{qualitative simulation}\index{compositional modeling}
